% SOURCES: docs/0.2.0v/architecture.md; docs/0.2.0v/element-model.md;
%          docs/1.0.0v/supervisor-report.md (§4–§6, §5a); README.md;
%          src/document/*, src/generator/*, src/export/*, src/draw/*, src/match/*, mcp/*

\chapter{Implementation}
\label{ch:impl}

This chapter describes how the system was built, following the data as it flows
from the Document Model through the generator to the export pipeline, and then the
three v1.0.0 authoring surfaces and the security boundary.

\section{The Document Model}
% SOURCES: src/document/types.ts, operations.ts; element-model.md
Every page is one \texttt{Document}: a tree of \texttt{ElementNode}s together with
the design tokens, SEO configuration, runtime flags, variables, settings, and asset
manifest (the class model in Figure~\ref{fig:class}). The primitive types ---
container, text, image, button, link, icon, list, and divider --- are the only node
kinds; complex shapes such as a hero or a cards grid are \emph{presets} that compose
primitives, not new classes. This ``composition over enumeration'' rule keeps the
model small and the generator simple.

All mutation flows through a single discriminated \texttt{Operation} union. Each
operation is a pure function over an \texttt{immer} draft --- signature
\texttt{(draft: Document, op: Op) => void} --- applying exactly one logical change,
and each carries a \texttt{kind} discriminator so the store can capture patches for
undo/redo. Listing~\ref{lst:ops} shows two representative operations; note that
style writes are targeted at a named breakpoint slot, so an edit to the tablet
value never silently overwrites the desktop base.

\begin{lstlisting}[language=,caption={Two operations from the discriminated union (\texttt{src/document/operations.ts}).},label=lst:ops]
/** Insert an element under an existing container. */
export interface InsertElementOp {
  readonly kind: 'insertElement'
  readonly parentId: ElementId
  readonly node: ElementNode
  /** Position among parent.children; appended when omitted. */
  readonly index?: number
}

/**
 * Write a value into node.style[breakpoint][...path]. The named
 * breakpoint slot is created if absent; siblings are left untouched.
 */
export interface UpdateNodeStyleOp {
  readonly kind: 'updateNodeStyle'
  readonly id: ElementId
  readonly breakpoint: BreakpointKey
  readonly path: PropertyPath
  readonly value: unknown
}
\end{lstlisting}

Files load through a strict boundary: a \texttt{.dtw} file is parsed, validated
with Zod, run through any version migrations, and validated again, so an old or
corrupt file can never reach the internal types unchecked.

\section{The Generator}
% SOURCES: src/generator/index.ts, htmlEmitter.ts, cssEmitter.ts; architecture.md §5
The generator is a compiler, not an inferer: it walks the document and emits
HTML, CSS, and JS deterministically, so identical input yields byte-identical
output. Tag selection is driven by each node's semantic role (with a sensible
default per type), which is how the output earns its semantic structure. The CSS is
tokens-driven --- a \texttt{:root} block defines every token, dark-theme overrides
and a \texttt{prefers-color-scheme} block follow, and element rules reference
\texttt{var(--token)} rather than raw values. Layout is expressed with CSS Grid,
Flexbox, and \texttt{clamp()}, and never with absolute positioning; this invariant
is enforced by a test that scans the generated CSS.

\section{Runtime Snippets}
% SOURCES: src/runtime/*, src/generator/jsEmitter.ts; supervisor-report §5
JavaScript is opt-in per behaviour: theme toggle, scroll-spy, smooth scroll, mobile
navigation, nav-on-scroll, reveals, animation gating, and terminal typing are each
gated by a flag that defaults to off. The JS emitter concatenates only the enabled
snippets into a single IIFE; if every flag is off, no \texttt{<script>} is emitted
at all. Every snippet is written to the same discipline --- passive event
listeners, work deferred with \texttt{requestAnimationFrame} or debounced, a
no-operation when its feature is disabled, and full respect for
\texttt{prefers-reduced-motion} --- so enabling a behaviour never degrades the
page's performance or accessibility.

\section{The Export Pipeline}
% SOURCES: src/export/index.ts; src/seo/axeGate.ts; supervisor-report §6
Export (Figures~\ref{fig:seqexport} and~\ref{fig:exportflow}) chains nine stages
--- \emph{validate, generate, inject-SEO, a11y-gate, optimize-images, minify,
sitemap/robots, bundle, save} --- each emitting a structured progress event so the
UI can show exactly where the pipeline is. The result is a discriminated union, so
a failure reports precisely which stage failed and why. The accessibility gate runs
\emph{before} any asset crosses the IPC boundary to the main process, so a failing
document never reaches image optimisation or packaging.

\section{v1.0.0 Authoring Surfaces}
% SOURCES: src/draw/*, src/match/*, mcp/*; supervisor-report §5a; README
\subsection{Draw-to-create}
A drawn rectangle is normalised, classified by \texttt{interpretRectangle}
(a ranked kind with a confidence score and a short explainer), and resolved to a
grid placement by \texttt{snapToGrid} (pure math over column start/span and an
insertion index --- never pixels). The canvas then builds an ordinary node with the
existing primitive factory and dispatches a normal \texttt{insertElement} operation,
so the drawn pixels never enter the store (Figure~\ref{fig:seqdraw}).

\subsection{Match-layout}
\texttt{extractSignature} builds a deterministic structural fingerprint of a
document tree (again, never pixels); \texttt{matchLayout} ranks that fingerprint
against six bundled professional pages --- each shipping a build-time precomputed
signature --- with a per-dimension breakdown; and the chosen page is hydrated
through the very same entry point that a project load uses. Adopting a design is
therefore indistinguishable, downstream, from opening a file.

\subsection{Headless MCP server}
A \texttt{stdio} Model Context Protocol server exposes roughly twenty tools over
the same pipeline (\texttt{create\_document}, \texttt{insert\_preset},
\texttt{apply\_template}, \texttt{set\_tokens}, \texttt{match\_layout},
\texttt{run\_a11y\_check}, \texttt{export\_site}, and so on). Each tool is a thin
adapter over the existing operations and entry points --- no logic is
reimplemented --- and it runs without Electron by using a filesystem-backed export
shim. Because it reuses the same pipeline, a headless export passes the same
accessibility gate as an interactive one.

\section{Process Boundaries and Security}
% SOURCES: src/main/*, src/preload/*; plan.md §15; supervisor-report §3
The main process runs with context isolation, no Node integration, and sandboxing;
the preload exposes only a typed \texttt{electronAPI} bridge, so the renderer can
never reach raw \texttt{ipcRenderer} or Node built-ins. Every IPC handler validates
its input --- path sanitisation, size caps, and MIME sniffing on uploads --- and a
strict Content Security Policy is applied in production (relaxed only for
hot-module reload in development). Validation happens at the edges; the internal
types are trusted in between.
